\subsection{Rank-Nullity and Invertibility}
\begin{proposition} \label{prop:injective-iff-kernel-zero} \hlt{Injectivity and the Kernel}\\
A linear map $T \in \LL(\mathcal{V}, \mathcal{W})$ is injective if and only if $\ker(T) = \{ \mathbf{0} \}$.
\begin{proof}
\begin{enumerate}[label=\roman*.]
\setlength{\itemsep}{0pt}
\item[$(\Rightarrow)$] If $T$ is injective and $\mathbf{v}\in\ker(T)$, then $T\mathbf{v}=\mathbf{0}=T\mathbf{0}$, so $\mathbf{v}=\mathbf{0}$.
\item[$(\Leftarrow)$] If $\ker(T)=\{\mathbf{0}\}$ and $T\mathbf{v}=T\mathbf{w}$, then $T(\mathbf{v}-\mathbf{w})=\mathbf{0}$. Hence $\mathbf{v}-\mathbf{w}\in\ker(T)$, so $\mathbf{v}=\mathbf{w}$.
\end{enumerate}
\hlt{Structural Commentary}: The kernel measures exactly where a linear map collapses vectors to zero.\\
\hlt{Key Insight}: Injectivity is equivalent to collapsing only the zero vector.
\end{proof}
\end{proposition}

\begin{theorem} \label{thm:range-spanned-by-images-of-basis} \hlt{Range from Basis Images}\\
Let $\mathcal{V}$ be finite-dimensional with basis $(\mathbf{v}_1,\ldots,\mathbf{v}_n)$, and let $T\in\LL(\mathcal{V},\mathcal{W})$. Then $\operatorname{range}(T)=\operatorname{span}(T\mathbf{v}_1,\ldots,T\mathbf{v}_n)$.\\
In particular, $\operatorname{range}(T)$ is finite-dimensional and $\dim(\operatorname{range}(T))\leq n$.
\begin{proof}
If $\mathbf{v}=c_1\mathbf{v}_1+\cdots+c_n\mathbf{v}_n$, then $T\mathbf{v}=c_1T\mathbf{v}_1+\cdots+c_nT\mathbf{v}_n$.\\
Thus every vector in $\operatorname{range}(T)$ lies in $\operatorname{span}(T\mathbf{v}_1,\ldots,T\mathbf{v}_n)$, and the reverse inclusion is immediate from the definition of range. The dimension bound follows because this spanning list has length $n$.\\
\hlt{Structural Commentary}: The range is generated by the images of any basis of the domain.\\
\hlt{Key Insight}: Linear maps carry a basis-spanning formula directly into a range-spanning formula.
\end{proof}
\end{theorem}

\begin{definition} \label{def:rank-and-nullity-linear-map} \hlt{Rank and Nullity}\\
For a finite-dimensional domain, the rank of $T$ is $\dim(\operatorname{range}(T))$ and the nullity of $T$ is $\dim(\ker(T))$.
\end{definition}

\begin{theorem} \label{thm:rank-nullity} \hlt{Rank-Nullity Theorem}\\
If $\mathcal{V}$ is finite-dim and $T \in \LL(\mathcal{V}, \mathcal{W})$, then $\operatorname{range}(T)$ is finite-dim and $\dim(\mathcal{V}) = \dim(\ker(T)) + \dim(\operatorname{range}(T))$.
\begin{proof}
By Proposition \ref{prop:kernel-and-range-are-subspaces} and Corollary \ref{cor:finite-dimensional-spaces-and-subspaces-have-bases}, $\ker(T)$ has a basis $(\mathbf{u}_1,\ldots,\mathbf{u}_m)$.\\
Extend it to a basis $(\mathbf{u}_1,\ldots,\mathbf{u}_m,\mathbf{w}_1,\ldots,\mathbf{w}_r)$ of $\mathcal{V}$ by Theorem \ref{thm:extending-independent-lists}.\\
If $\mathbf{v}=\sum\limits_{i=1}^m a_i\mathbf{u}_i+\sum\limits_{j=1}^r b_j\mathbf{w}_j$, then $T\mathbf{v}=\sum\limits_{j=1}^r b_jT\mathbf{w}_j$, because $T\mathbf{u}_i=\mathbf{0}$. Hence $(T\mathbf{w}_1,\ldots,T\mathbf{w}_r)$ spans $\operatorname{range}(T)$.\\
If $\sum\limits_{j=1}^r c_jT\mathbf{w}_j=\mathbf{0}$, then $\sum\limits_{j=1}^r c_j\mathbf{w}_j\in\ker(T)$, so for some scalars $d_i$, $\sum\limits_{j=1}^r c_j\mathbf{w}_j=\sum\limits_{i=1}^m d_i\mathbf{u}_i$.\\
Thus $\sum\limits_{j=1}^r c_j\mathbf{w}_j-\sum\limits_{i=1}^m d_i\mathbf{u}_i=\mathbf{0}$, and linear independence of extended basis forces every $c_j$ and $d_i$ to be zero.\\
Hence $(T\mathbf{w}_1,\ldots,T\mathbf{w}_r)$ is a basis of $\operatorname{range}(T)$. Thus $\dim(\ker(T))=m$, $\dim(\operatorname{range}(T))=r$, and $\dim(\mathcal{V})=m+r$.\\
\hlt{Structural Commentary}: An adapted basis splits $\mathcal{V}$ into kernel directions and directions detected by $T$.\\
\hlt{Key Insight}: Rank-nullity says every domain dimension is either killed by $T$ or appears in its range.
\end{proof}
\end{theorem}

\begin{corollary} \label{cor:map-to-smaller-space-not-injective} \hlt{Maps to Smaller Spaces Are Not Injective}\\
Let $\mathcal{V}$ and $\mathcal{W}$ be finite-dimensional vector spaces with $\dim(\mathcal{V})>\dim(\mathcal{W})$.\\
Then no linear map from $\mathcal{V}$ to $\mathcal{W}$ is injective.
\begin{proof}
Let $T\in\LL(\mathcal{V},\mathcal{W})$. By Theorem \ref{thm:rank-nullity}, $\dim(\ker(T))=\dim(\mathcal{V})-\dim(\operatorname{range}(T)) \geq \dim(\mathcal{V})-\dim(\mathcal{W})>0$.\\
Thus $\ker(T)$ contains a nonzero vector, so $T$ is not injective by Proposition \ref{prop:injective-iff-kernel-zero}.\\
\hlt{Structural Commentary}: Excess domain dimension forces collapse.\\
\hlt{Key Insight}: The dimension deficit must appear inside the kernel.
\end{proof}
\end{corollary}

\begin{corollary} \label{cor:map-to-larger-space-not-surjective} \hlt{Maps to Larger Spaces Are Not Surjective}\\
Let $\mathcal{V}$ and $\mathcal{W}$ be finite-dimensional vector spaces with $\dim(\mathcal{V})<\dim(\mathcal{W})$.\\
Then no linear map from $\mathcal{V}$ to $\mathcal{W}$ is surjective.
\begin{proof}
Let $T\in\LL(\mathcal{V},\mathcal{W})$. By Theorem \ref{thm:rank-nullity}, $\dim(\operatorname{range}(T))=\dim(\mathcal{V})-\dim(\ker(T))\leq\dim(\mathcal{V})<\dim(\mathcal{W})$.\\
Hence $\operatorname{range}(T)$ is a proper subspace of $\mathcal{W}$, so $T$ is not surjective.\\
\hlt{Structural Commentary}: A smaller domain cannot cover a larger codomain.\\
\hlt{Key Insight}: The range can never have dimension larger than the domain.
\end{proof}
\end{corollary}

\begin{corollary} \label{cor:equal-dimension-kernel-range} \hlt{Equal-Dimension Kernel-Range Criterion}\\
Let $\mathcal{V}$ and $\mathcal{W}$ be finite-dimensional vector spaces with $\dim(\mathcal{V})=\dim(\mathcal{W})$, and let $T\in\LL(\mathcal{V},\mathcal{W})$.\\
Then $\ker(T)=\{\mathbf{0}\}$ if and only if $\operatorname{range}(T)=\mathcal{W}$. Equivalently, $T$ is injective if and only if $T$ is surjective.
\begin{proof}
\begin{enumerate}[label=\roman*.]
\setlength{\itemsep}{0pt}
\item[$(\Rightarrow)$] If $\ker(T)=\{\mathbf{0}\}$, then rank-nullity gives $\dim(\operatorname{range}(T))=\dim(\mathcal{V})=\dim(\mathcal{W})$.\\
Since $\operatorname{range}(T)$ is a subspace of $\mathcal{W}$ with the same dimension, $\operatorname{range}(T)=\mathcal{W}$.
\item[$(\Leftarrow)$] If $\operatorname{range}(T)=\mathcal{W}$, then rank-nullity gives $\dim(\ker(T))=\dim(\mathcal{V})-\dim(\mathcal{W})=0$, hence $\ker(T)=\{\mathbf{0}\}$.
\end{enumerate}
\hlt{Structural Commentary}: Equal finite dimensions make injectivity and surjectivity equivalent.\\
\hlt{Key Insight}: In equal dimensions, having no kernel loss is exactly having full range.
\end{proof}
\end{corollary}

\begin{definition} \hlt{Invertible Linear Map and Isomorphism}\\
A linear map $T \in \LL(\mathcal{V}, \mathcal{W})$ is \hlt{invertible} if there exists $S \in \LL(\mathcal{W}, \mathcal{V})$ such that $ST = I_{\mathcal{V}}$ and $TS = I_{\mathcal{W}}$.\\
An invertible linear map is called an \hlt{isomorphism}, and if one exists we write $\mathcal{V} \cong \mathcal{W}$.
\end{definition}

\begin{proposition} \label{prop:inverse-linear-map-is-unique} \hlt{Inverse Is Unique}\\
An invertible linear map has a unique inverse. Thus if $T$ is invertible, its inverse is denoted $T^{-1}$.
\begin{proof}
Suppose $S_1,S_2\in\LL(\mathcal{W},\mathcal{V})$ are inverses of $T\in\LL(\mathcal{V},\mathcal{W})$.\\
Then $S_1 = S_1I_{\mathcal{W}} = S_1(TS_2) = (S_1T)S_2 = I_{\mathcal{V}}S_2 = S_2$. Hence the inverse is unique.\\
\hlt{Structural Commentary}: A two-sided inverse is forced by the composition identities.\\
\hlt{Key Insight}: Associativity lets the two candidate inverses meet around $T$ and collapse to the same map.
\end{proof}
\end{proposition}

\begin{proposition} \label{prop:inverse-of-product-linear-maps} \hlt{Inverse of a Product Reverses Order}\\
Let $T\in\LL(\mathcal{U},\mathcal{V})$ and $S\in\LL(\mathcal{V},\mathcal{W})$ be invertible. Then $ST$ is invertible and $(ST)^{-1}=T^{-1}S^{-1}$.
\begin{proof}
Since $T^{-1}\in\LL(\mathcal{V},\mathcal{U})$ and $S^{-1}\in\LL(\mathcal{W},\mathcal{V})$, the product $T^{-1}S^{-1}\in\LL(\mathcal{W},\mathcal{U})$ is defined.\\
Moreover, $(T^{-1}S^{-1})(ST)=T^{-1}(S^{-1}S)T=T^{-1}T=I_{\mathcal{U}}$ and $(ST)(T^{-1}S^{-1})=S(TT^{-1})S^{-1}=SS^{-1}=I_{\mathcal{W}}$.\\
Thus $T^{-1}S^{-1}$ is a two-sided inverse of $ST$, so $(ST)^{-1}=T^{-1}S^{-1}$.\\
\hlt{Structural Commentary}: To undo a composition, undo the last map first.\\
\hlt{Key Insight}: Composition order reverses under inversion.
\end{proof}
\end{proposition}

\begin{proposition} \label{prop:linear-map-invertible-iff-bijective} \hlt{Invertibility and Bijectivity}\\
A linear map is invertible if and only if it is both injective and surjective.
\begin{proof}
\begin{enumerate}[label=\roman*.]
\setlength{\itemsep}{0pt}
\item[$(\Rightarrow)$] Let $S$ be an inverse of $T$. If $T\mathbf{v}=T\mathbf{w}$, then $\mathbf{v}=ST\mathbf{v}=ST\mathbf{w}=\mathbf{w}$, so $T$ is injective. Also every $\mathbf{w}$ in the codomain satisfies $\mathbf{w}=T(S\mathbf{w})$, so $T$ is surjective.
\item[$(\Leftarrow)$] Suppose $T$ is injective and surjective. Define $S\mathbf{w}$ to be the unique vector $\mathbf{v}$ with $T\mathbf{v}=\mathbf{w}$.\\
Since both sides below map under $T$ to $a\mathbf{w}_1+\mathbf{w}_2$, $S(a\mathbf{w}_1+\mathbf{w}_2)=aS\mathbf{w}_1+S\mathbf{w}_2$. Thus $S$ is linear and satisfies $ST=I$ and $TS=I$, so $T$ is invertible.
\end{enumerate}
\hlt{Structural Commentary}: Bijectivity gives a unique reverse map, and linearity makes it a linear inverse.\\
\hlt{Key Insight}: A linear map is invertible exactly when it is one-to-one and onto.
\end{proof}
\end{proposition}

\begin{definition} \label{def:linear-map-morphism-terminology-vector-spaces} \hlt{Morphism Terminology for Linear Maps}\\
For $T\in\LL(\mathcal{V},\mathcal{W})$:
\begin{enumerate}[label=(\roman*)]
\setlength{\itemsep}{0pt}
\item $T$ is a \hlt{homomorphism}.
\item If $\mathcal{V}=\mathcal{W}$, then $T$ is an \hlt{endomorphism}.
\item If $\ker(T)=\{\mathbf{0}\}$, then $T$ is a \hlt{monomorphism}, or \hlt{embedding}.
\item If $\operatorname{range}(T)=\mathcal{W}$, then $T$ is an \hlt{epimorphism}.
\item If $\ker(T)=\{\mathbf{0}\}$ and $\operatorname{range}(T)=\mathcal{W}$, then $T$ is an \hlt{isomorphism}.
\item If $\mathcal{V}=\mathcal{W}$ and $T$ is invertible, then $T$ is an \hlt{automorphism}.
\end{enumerate}
\end{definition}

\begin{proposition} \label{prop:rank-nullity-complement-view} \hlt{Rank-Nullity as a Complement Isomorphism}\\
Let $T\in\LL(\mathcal{V},\mathcal{W})$, and suppose $\mathcal{V}=\ker(T)\oplus\mathcal{M}$. Then $T|_{\mathcal{M}}:\mathcal{M}\to\operatorname{range}(T)$ is an isomorphism.
\begin{proof}
Indeed, $\ker(T|_{\mathcal{M}})=\mathcal{M}\cap\ker(T)=\{\mathbf{0}\}$, and every $T\mathbf{v}$ equals $T\mathbf{m}$ for the decomposition $\mathbf{v}=\mathbf{k}+\mathbf{m}$.\\
Thus rank-nullity counts $\mathcal{V}=\ker(T)\oplus\mathcal{M}$ with $\mathcal{M}\cong\operatorname{range}(T)$.
\end{proof}
\end{proposition}
